% Generado por scripts/70_tablas_figuras.R a partir de resultados-v1. No editar a mano.
\begin{table}[htbp]
\centering
\small
\caption{Tabla 6.9—Efecto mínimo detectable de la evaluación por polígonos}
\label{tab:tabla_6_9}
\begin{threeparttable}
\begin{tabular}{@{}lcccccc@{}}
\toprule
 & (1) & (2) & (3) & (4) & (5) & (6) \\
 & Proyectos & Celdas & Media previa & ATT & EMD & EMD / media \\
\midrule
2013 & 7 & 4.248 & 0,53 & $-$0,24 & 1,38 & 258 \% \\
 &  &  &  & (0,49) &  &  \\
2014 & 4 & 718 & 1,46 & 0,13 & 3,75 & 257 \% \\
 &  &  &  & (1,34) &  &  \\
2015 & 4 & 919 & 2,02 & 0,01 & 5,72 & 283 \% \\
 &  &  &  & (2,04) &  &  \\
2017 & 3 & 1.290 & 2,34 & 0,41 & 5,34 & 228 \% \\
 &  &  &  & (1,90) &  &  \\
2018 & 7 & 9.889 & 0,12 & $-$0,11 & 0,41 & 338 \% \\
 &  &  &  & (0,15) &  &  \\
\midrule
Todas las cohortes & 25 & 17.064 & 0,52 & $-$0,10 & 0,98 & 190 \% \\
 &  &  &  & (0,35) &  &  \\
\bottomrule
\end{tabular}
\begin{tablenotes}[flushleft]
\footnotesize
\item \textit{Notas:} Pérdida de bosque natural en hectáreas por celda de 400 ha y año. Columna (4): ATT con su error estándar entre paréntesis. EMD: efecto mínimo detectable con 5 \% de significancia (dos colas) y 80 \% de potencia, igual a 2,8 veces el error estándar. Columna (6): EMD como porcentaje de la pérdida previa de las celdas tratadas; un valor mayor que 100 \% significa que ni siquiera eliminar toda la deforestación sería detectable. No hay cohorte 2016 ni 2019 con polígono.
\item \textit{Fuente:} MapBiomas Colombia, Colección 3; Karnik et al. (2025); cálculos propios.
\end{tablenotes}
\end{threeparttable}
\end{table}
